\chapter{Arithmetic \(E_{n}\) Formality and CoHA}
\label{v4-arith:kontsevich-soibelman}

\section{\(E_{1}\)-Formality Is Automatic over \(\mathbb{Z}\)}
\label{v4-arith:ks:E1-automatic}

\begin{theorem}[Fresse \emph{Homotopy of Operads}, Vol.~II, Thm~I.4.1.6]
\label{v4-arith:thm:E1-formality-automatic}
\ClaimStatusProvedElsewhere. For any commutative ring \(R\),
\(H_{\bullet}(E_{1};R)\cong\mathrm{Ass}_{R}\) and the canonical map
\(C_{\bullet}(E_{1};R)\to\mathrm{Ass}_{R}\) is an operadic
quasi-isomorphism. \(E_{1}\)-formality is automatic over \(\mathbb{Z}\).
\end{theorem}

\begin{corollary}[ordered \(E_{1}\)-chiral survives over \(\mathbb{Z}\)]
\label{v4-arith:cor:ordered-E1-survives}
\ClaimStatusProvedHere. The ordered \(E_{1}\)-chiral structure on
\(\mathcal{V}^{\mathrm{prim}}\), whose underlying product is the
strict associative product from the ordered factorization-algebra
product, is defined integrally: \(H_{\bullet}(E_{1};\mathbb{Z})\cong\mathrm{Ass}_{\mathbb{Z}}\)
carries no torsion, so no Massey obstruction arises.
\end{corollary}

\section{\(E_{2}\)-Formality Fails Integrally}
\label{v4-arith:ks:E2-fails}

\begin{theorem}[Salvatore--Wahl; Cohen LNM~533; Fresse--Willwacher arXiv:1503.08699]
\label{v4-arith:thm:E2-fails-integrally}
\ClaimStatusProvedElsewhere. \(H_{\bullet}(E_{2};\mathbb{F}_{p})\)
carries non-trivial Dyer--Lashof operations; these obstruct
formality because they vanish on any cohomology-operadic target
but survive on the chain-level operad. Over \(\mathbb{Q}\), Tamarkin's
theorem (arXiv:math/9803025) establishes
\(C_{\bullet}(E_{2};\mathbb{Q})\xrightarrow{\sim}H_{\bullet}(E_{2};\mathbb{Q})=\mathrm{Ger}_{\mathbb{Q}}\)
(where \(\mathrm{Ger}_{\mathbb{Q}}\) is the Gerstenhaber operad over \(\mathbb{Q}\)),
a torsor for the Grothendieck--Teichm\"{u}ller group
\(\mathrm{GRT}_{1}(\mathbb{Q})\).
\end{theorem}

\begin{corollary}[no integral \(E_{2}\)-enhancement]
\label{v4-arith:cor:no-integral-E2}
\ClaimStatusProvedHere. The \(E_{2}\)-chiral-centre enhancement of
\(\mathcal{V}^{\mathrm{prim}}\) (Drinfeld-centre or Swiss-cheese)
does not extend to \(\mathrm{Spec}\,\mathbb{Z}\). It lifts to
\(\mathbb{Q}\) as a \(\mathrm{GRT}_{1}(\mathbb{Q})\)-torsor.
The lifting to \(\mathbb{Z}[1/N]\), for \(N\) the exponent of the
torsion of \(H_{\bullet}(E_{2};\mathbb{Z})\), is expected from
Cohen's computation (LNM~533) but has not been verified in this text;
it is recorded as heuristic pending that computation.
\end{corollary}

\begin{remark}[Kaledin Japanese cocycle]
Kaledin isolates the arithmetic obstruction: for
an associative \(\mathbb{Z}\)-algebra \(A\), the reduction
\(A\otimes\mathbb{F}_{p}\) may carry a non-trivial Hochschild
\(A_{\infty}\)-deformation class, the ``Japanese cocycle,'' in
\(\mathrm{HH}^{2}(A\otimes\mathbb{F}_{p})/\mathrm{HH}^{2}(A)\otimes\mathbb{F}_{p}\).
This class is the first-order obstruction to integral strictification of
the \(E_{2}\)-enhancement.
\end{remark}

\section{Higher Products Vanish}
\label{v4-arith:ks:higher-products-vanish}

\(E_1\)-formality identifies \(C_{\bullet}(E_1;\mathbb{Z})\) with
\(\mathrm{Ass}_{\mathbb{Z}}\) as operads; \(\mathrm{Ass}_{\mathbb{Z}}\)
has only \(m_1\) and \(m_2\) as non-trivial structure maps.
The question is whether this strictness persists for
\(\mathcal{V}^{\mathrm{prim}}\): do the \(A_{\infty}\)-higher products
\(m_k\) for \(k \geq 3\) vanish?

\begin{theorem}
\label{v4-arith:thm:mk-vanish}
\ClaimStatusProvedHere. The \(A_{\infty}\)-higher products \(m_{k}\)
(\(k\geq 3\)) vanish on \(\mathcal{V}^{\mathrm{prim}}\) over
\(\mathbb{Z}\), because \(H_{\bullet}(E_{1};\mathbb{Z})\cong\mathrm{Ass}_{\mathbb{Z}}\)
is torsion-free and carries no higher operations.
\end{theorem}

\begin{corollary}[strict associativity over \(\mathbb{Z}\)]
\label{v4-arith:cor:strict-associativity-Z}
\ClaimStatusProvedHere. \(\mathcal{V}^{\mathrm{prim}}\) is strictly
associative over \(\mathbb{Z}\): the \(A_{\infty}\)-algebra structure
is equivalent to a strict dg-algebra with no higher homotopies.
The \(E_{2}\)-chiral-centre obstructions of
\S\ref{v4-arith:ks:E2-fails} operate at a separate level and are
unaffected by this strict associativity.
\end{corollary}

\section{The Chiral Bracket via Adelic Residue}
\label{v4-arith:ks:chiral-bracket-adelic}

\begin{definition}[arithmetic \(\lambda\)-bracket]
\label{v4-arith:def:arith-lambda-bracket}
Let \(\overline{C}\) denote the Arakelov compactification of the
arithmetic curve, obtained by adjoining the archimedean place
\(|\cdot|_{\infty}\) to \(\mathrm{Spec}\,\mathbb{Z}\).
The chiral bracket on \(\mathcal{V}^{\mathrm{prim}}\) over
\(\overline{C}\) is the Beilinson--Drinfeld \(\lambda\)-bracket on
\(\mathrm{Conf}_{2}(\overline{C})=\overline{C}\times\overline{C}\setminus\Delta\):
\begin{equation}
\{a,b\}_{\mathrm{arith}}(p,q)\;=\;\mathrm{Res}_{p=q}\,\omega_{a,b}(p,q)\cdot\log|p-q|_{\infty}^{-1},
\end{equation}
where \(\omega_{a,b}\) is the section of
\(\mathcal{V}^{\mathrm{prim}}\boxtimes\mathcal{V}^{\mathrm{prim}}\)
over \(\mathrm{Conf}_{2}(\overline{C})\) determined by \(a\) and \(b\).
Each finite prime contributes a discrete residue at the corresponding
closed point; the archimedean place contributes the
\(|\cdot|_{\infty}\) normalisation.
\end{definition}

\begin{theorem}[chiral bracket descends adelically]
\label{v4-arith:thm:chiral-bracket-adelic}
\ClaimStatusConjectured. The arithmetic \(\lambda\)-bracket on
\(\mathcal{V}^{\mathrm{prim}}\) is a well-defined adelic residue,
interpreted via Arakelov completion: the arithmetic cotangent bundle
\(\mathrm{Cot}^{*}_{\mathrm{arith}}\) (the sheaf of arithmetic
\(1\)-forms on \(\overline{C}\), metrized at the archimedean place)
on \(\mathrm{Conf}_{2}(\overline{C})\) restricted to the punctured
diagonal supports the pairing
\(\{\mathcal{V}^{\mathrm{prim}},\mathcal{V}^{\mathrm{prim}}\}_{\mathrm{arith}}\subset\mathcal{V}^{\mathrm{prim}}\otimes_{\mathbb{Z}}\mathrm{Cot}^{*}_{\mathrm{arith}}\).
\end{theorem}

\section{The Arithmetic CoHA Conjecture}
\label{v4-arith:ks:CoHA-conjectural}

Given that \(\mathcal{V}^{\mathrm{prim}}\) carries an integral
\(E_1\)-structure and an adelic chiral bracket, one asks whether
it fits into the Kontsevich--Soibelman cohomological Hall algebra
(CoHA) framework as a BPS sub-object.

\begin{theorem}[Kontsevich--Soibelman arXiv:1006.2706]
\label{v4-arith:thm:KS-CoHA}
\ClaimStatusProvedElsewhere. For a quiver \(Q\) and potential
\(W\in\mathbb{C}\langle Q\rangle_{\mathrm{cyc}}\), the cohomological
Hall algebra is
\begin{equation*}
\mathrm{CoHA}(Q,W)\;:=\;\bigoplus_{\mathbf{d}}H^{\bullet}_{G_{\mathbf{d}}}(\mathrm{Rep}(Q,\mathbf{d}),\varphi_{W_{\mathbf{d}}}),
\end{equation*}
where \(\mathbf{d}\) ranges over dimension vectors,
\(G_{\mathbf{d}}=\prod_{v}\mathrm{GL}_{d_v}\) is the gauge group,
\(\mathrm{Rep}(Q,\mathbf{d})\) the representation variety, and
\(\varphi_{W_{\mathbf{d}}}\) the vanishing cycle sheaf of the
trace of \(W\) restricted to \(\mathrm{Rep}(Q,\mathbf{d})\).
\end{theorem}

\begin{theorem}[Davison integrality arXiv:1601.02479]
\label{v4-arith:thm:davison-integrality}
\ClaimStatusProvedElsewhere. For symmetric \((Q,W)\) in the CY-3
setting,
\begin{equation*}
\mathrm{CoHA}(Q,W)\;\cong\;\mathrm{Sym}(\mathrm{BPS}(Q,W)\otimes H^{\bullet}(B\mathbb{C}^{\times})),
\end{equation*}
with \(\mathrm{BPS}(Q,W)\) the BPS Lie algebra (primitive generator
lattice, in Davison's sense).
\end{theorem}

\begin{theorem}[conjectural arithmetic CoHA identification]
\label{v4-arith:thm:arithmetic-CoHA-conjectural}
\ClaimStatusConjectured. \(\mathcal{V}^{\mathrm{prim}}\) is the
\emph{BPS sub-object}
\begin{equation*}
\mathcal{V}^{\mathrm{prim}}\;\stackrel{?}{\cong}\;\mathrm{BPS}(Q^{\mathrm{arith}},W^{\mathrm{arith}})
\end{equation*}
of the full \(\mathrm{CoHA}(Q^{\mathrm{arith}},W^{\mathrm{arith}})\),
where the full CoHA is a polynomial algebra on \(\mathcal{V}^{\mathrm{prim}}\)
via Davison's theorem (\Cref{v4-arith:thm:davison-integrality}).
The candidate arithmetic quiver-with-potential is:
\begin{itemize}
\item \(Q^{\mathrm{arith}}_{0}=\mathbb{P}\cup\{\infty\}\) (primes plus
the archimedean place);
\item \(Q^{\mathrm{arith}}_{1}=\) Hecke correspondences (one edge per
automorphism of a local completion);
\item \(W^{\mathrm{arith}}\in\mathrm{HC}_{0}(\mathbb{Z}\langle Q^{\mathrm{arith}}\rangle)\),
the cyclic cohomology class of the Connes--Consani motivic lift
of the global \(\zeta\)-function (Connes--Consani 2009,
\emph{Schemes over \(\mathbb{F}_{1}\) and zeta functions},
Compos.~Math.~146, arXiv:1006.4810, \S5; distinct from the
arithmetic-site arXiv:1405.4527 used elsewhere in this branch).
\end{itemize}
\end{theorem}

\begin{remark}[Bost--Connes quiver and the role of the potential]
The Bost--Connes quiver (vertex set \(\mathbb{P}\), free commutative
monoid on primes) is acyclic, with trivial potential; its CoHA is the
polynomial algebra \(\bigotimes_{p}\mathbb{Z}[e_{p}]\).
The identification of \(\mathcal{V}^{\mathrm{prim}}\) with
\(\mathrm{BPS}(Q^{\mathrm{arith}},W^{\mathrm{arith}})\) requires the
cyclic-cohomology-lifted potential \(W^{\mathrm{arith}}\), which encodes
the \(\zeta\)-function data encoded in \(\mathrm{HC}_0\).
\end{remark}

\section{\(\log L\) via Cyclic Cohomology}
\label{v4-arith:ks:logL-not-cyclic}

\begin{observation}
\label{v4-arith:obs:logL-not-cyclic}
\(\log L=\sum_{p}\sum_{k\geq 1}p^{-ks}/k\) is a formal Dirichlet
series with transcendental coefficients \(\log p\); its image in the
path algebra \(\mathbb{Z}\langle Q\rangle_{\mathrm{cyc}}\) is defined
only after passage through Connes' cyclic cohomology \(\mathrm{HC}_{0}\),
which accommodates the \(\log p\) coefficients as cyclic cohomology
classes rather than as algebraic potential terms.
\end{observation}

\begin{remark}[role of exponential Hodge structure]
Kontsevich--Soibelman \emph{Cohomological Hall algebra,
exponential Hodge structures, and motivic DT invariants}
(arXiv:1006.2706, \S7) provides the setting in which \(\zeta\)-value
periods appear as CoHA data. The arithmetic CoHA identification
requires this exponential Hodge structure machinery.
\end{remark}

\section{Arithmetic Formality Stratification}
\label{v4-arith:ks:stratification}

The \(E_{\bullet}\)-formality of
\(\mathcal{V}^{\mathrm{prim}}\) over \(\overline{\mathrm{Spec}\,\mathbb{Z}}\)
stratifies by level: \(E_1\)-formality is integral
(\Cref{v4-arith:thm:E1-formality-automatic}), while
\(E_2\)-formality is obstructed integrally by Dyer--Lashof operations
and the Kaledin Japanese cocycle, and holds only over \(\mathbb{Q}\) as
a \(\mathrm{GRT}_1(\mathbb{Q})\)-torsor
(\Cref{v4-arith:cor:no-integral-E2}).
The strict associativity of \(\mathcal{V}^{\mathrm{prim}}\) over
\(\mathbb{Z}\) (\Cref{v4-arith:cor:strict-associativity-Z}) and the
adelic residue formula for the chiral bracket
(\Cref{v4-arith:def:arith-lambda-bracket}) are consequences of
\(E_1\)-integrality.
The identification of \(\mathcal{V}^{\mathrm{prim}}\) with the BPS
sub-object \(\mathrm{BPS}(Q^{\mathrm{arith}},W^{\mathrm{arith}})\)
of a conjectural arithmetic quiver-with-potential remains open:
the candidate \(W^{\mathrm{arith}}\) requires the cyclic cohomology
lift of \(\log L\) via exponential Hodge structure
(\Cref{v4-arith:thm:arithmetic-CoHA-conjectural}), and the quiver
\(Q^{\mathrm{arith}}\) itself has not been constructed.
